\section{Stage-Level Artifact Assessment}

The internal validation question is not simply whether the code executed. The
pipeline is useful only if each generated artifact is strong enough to support
the next stage. We therefore assess the reported Statute-of-Frauds run as a
source-to-score case study grounded in a real court case, \emph{Anglemire v.
Policemen's Benevolent Association of Chicago}. A single source case is passed through Frank,
Karthic, natural model-response generation, Dasha clustering, rubric-conditioned
judging, Zak escalation logic, and paper export. The question at each boundary
is whether the artifact is legally coherent, source-grounded, consistent with
the instruction examples, and auditable by a research team.

\begin{table}[h]
\centering
\caption{Stage-level artifact assessment for the live SOF case study.}
\label{tab:stage-artifact-assessment}
\begin{tabular}{p{0.13\linewidth}p{0.25\linewidth}p{0.32\linewidth}p{0.20\linewidth}}
\toprule
Stage & Validation question & Current SOF evidence & Internal assessment \\
\midrule
Frank & Is the source converted into a legally sound gold packet and neutral benchmark question? & Identified SOF marriage-provision and fraternal-benefit certificate issues; generated one neutral question, structured gold answer, and 3 variations. & Coherent source-grounded packet for the case; requires external legal review before publication. \\
Karthic & Is the rubric fresh, comprehensive, source-supported, and derived from Frank? & Produced 10 rows across doctrine, rule, facts, exceptions, counterargument, conclusion, variation, and source-support categories. & Good scoring core for internal validation; full row-card materialization remains a publication-hardening task. \\
Dasha & Are natural model answers clustered by legal reasoning, not labels or surface form? & Clustered 60 unlabeled natural responses into 23 track-aware reasoning clusters with all responses assigned. & Meaningful perturbation-aware separation; needs held-out legal cases before broad validity claims. \\
Judge & Is the rubric applied row by row to centroids and projected to members? & Produced 230 row-level scores, 60 projected member scores, and 10 model ranking entries. & Mechanism works; ranking is not yet a population-level model-performance claim. \\
Zak & Are failures and uncertainty made explicit? & The live SOF run emitted 1 Zak packet after judge disagreement/adjudication checks. & Packet creation works in the active live run; live threshold calibration remains. \\
\bottomrule
\end{tabular}

\end{table}

For Frank, the central question is whether the source case has been converted
into a legally sound benchmark rather than merely summarized. In the reported
run, Frank identified the Statute-of-Frauds marriage-provision theory and its
interaction with fraternal-benefit certificate rules; generated a neutral
Illinois-law question; produced a structured gold answer; and created three
boundary variations. In the perturbation lane, those variations become
executable question tracks rather than prose-only notes. The artifact is useful for technical review because it preserves the
legally operative facts: the premarital promise, marriage consideration,
certificate naming the wife, later replacement certificate, and association
rules. The remaining limitation is that the legal assessment is still internal:
the evidence supports the claim that Frank produced a coherent source-grounded
packet, not that outside experts have certified the gold answer.

For Karthic, the question is whether the rubric is a fresh scoring instrument
derived from Frank rather than a generic grading checklist. The reported run
produced 10 rows covering doctrine, rule, facts, exceptions, counterargument,
conclusion, variation sensitivity, and source support. That is the right shape
for the reported case because the answer space turns on doctrinal gate selection,
writing or certificate significance, later-beneficiary defenses, equitable
theories, and disciplined conclusion language. The reference implementation
stores the scoring core of each row rather than every full row-card field
specified in the ideal Karthic instructions. This is acceptable for the present
reported run, but the manuscript marks it as a limitation before publication.

For Dasha, the question is whether unlabeled natural model answers are grouped
by legal reasoning. The reported run used 60 natural responses across the
original question, an invariant party-name edit, and a material signed-writing
edit. It observed 15 track-aware clusters. The dominant original and invariant
families rejected the unwritten marriage-consideration promise and treated the
later bylaw-compliant designation as controlling; minority original clusters
captured equitable override, constructive-trust, and certificate-as-writing
theories. The material signed-writing track shifted the dominant answer family
toward written compliance and enforceability. This is the kind of separation
the pipeline is intended to find, because the clusters differ in legal theory
rather than merely in wording. For perturbation runs, Dasha clusters each
question track separately before comparing dominant signals, which prevents a
variant answer from being absorbed into the original question's cluster. The strongest
remaining risk is extraction error: Dasha in the reported configuration uses an LLM to extract a
reasoning signature, then groups exact normalized buckets. The method is
auditable, but future validation should repeat extraction, compare alternate
clusterers, and increase the number of natural model responses.

For Judge, the question is whether the rubric is actually applied row by row to
the Dasha centroids. The reported run generated 150 panel row scores: 10 rubric
rows across 15 clusters judged by a two-model panel. Six clusters with unstable
panel rows received an adjudication pass before final projection. The judge then
projected centroid scores to all 60 member responses and produced model
rankings. This demonstrates the intended
compression mechanism: score the legal-reasoning centroid, then propagate the
score to responses using that reasoning path. The present ranking is not a
population-level claim about model quality because the live batch is bounded
and uses one source case. It is evidence that the judge machinery can apply the
rubric, adjudicate unstable panel rows, and produce auditable rankings from
centroid scores.

For Zak, the validation question is whether uncertainty is represented
explicitly. In this run, the judge output required escalation after centroid
scoring and adjudication, and Zak emitted a disputed-cluster packet rather than
silently treating the ranking as final. The implementation also includes
positive offline regressions: one forces a low-margin centroid comparison to
create a disputed-cluster packet, and one forces unstable repeated judge rows to
create an unstable-row packet. These regressions show that Zak packet creation
is not merely an empty no-escalation record. The remaining work is to calibrate
the live thresholds across held-out cases so escalation is neither too rare nor
too noisy.

\subsection{Transfer Beyond Statute of Frauds}

The pipeline is designed to be doctrine-general, but the evidence in this
manuscript is intentionally narrower. The doctrine-transfer claim is therefore
architectural rather than empirical: the same source-to-packet schema,
instruction-context mechanism, natural-response protocol, reasoning-signature
extraction, centroid judging, and escalation logic can be reused for another
area of law if that area supplies adequate source material and role-specific
instruction context. Statute of Frauds is the calibration domain because it
contains multiple gates, exceptions, boundary variations, and plausible model
failure modes. Those properties make it a useful stressor for the general
pipeline design.

\begin{table}[h]
\centering
\caption{Doctrine-transfer argument and remaining validation burden.}
\label{tab:doctrine-transfer}
\begin{tabular}{p{0.24\linewidth}p{0.34\linewidth}p{0.32\linewidth}}
\toprule
Pipeline component & Doctrine-general mechanism & Remaining validation burden \\
\midrule
Frank & Uses source text plus instruction context to infer doctrine gates and create a packet schema that is not inherently SOF-only. An offline contract-interpretation smoke fixture now emits generic doctrine gates without a SOF packet. & Run live held-out non-SOF cases and review doctrine detection, question neutrality, and gold-answer quality. \\
Karthic & Builds rows from the locked Frank packet, so rubric content should follow the source doctrine rather than a fixed SOF checklist. The contract fixture produces a generic source-grounded rubric. & Confirm row coverage and source support in additional doctrine families. \\
Dasha & Extracts normalized reasoning signatures from natural answers after generation and can use Frank's source-derived gate aliases. The contract fixture separates buyer-remedy, contra-proferentem, and seller-limitation reasoning. & Test whether signature fields and buckets remain discriminative outside SOF on live natural model answers. \\
Judge & Applies row-level rubric criteria to centroids, independent of the doctrine label. & Repeat judging with alternate judge models, repeats, and expert comparison. \\
Zak & Escalates low agreement, mixed clusters, or uncertain source support using artifact metadata. & Add positive failure cases across domains to verify escalation sensitivity. \\
\bottomrule
\end{tabular}

\end{table}

The current implementation also includes an offline contract-interpretation
transfer smoke test. That fixture verifies that the artifact schema and Dasha
normalization path can run outside Statute of Frauds, but it is intentionally
reported as schema-transfer evidence rather than broad legal-domain validity.
The strongest future claim would require held-out live cases from other legal
domains. For each domain, the same stage-level questions should be asked: did
Frank infer the correct doctrinal elements from the source; did Karthic build a
useful rubric from that packet; did Dasha recover and cluster natural model
reasoning paths; did Judge apply the rubric consistently; and did Zak surface
failures instead of hiding them? Until that evidence exists, the paper should
claim that the method is designed for doctrine transfer, internally demonstrated
on Statute of Frauds, and smoke-tested on one non-SOF fixture, not that broad
legal-domain validity has already been proven.
